\documentclass[a4paper, final]{article}

\usepackage[russian]{babel}
\usepackage{fontspec}
\setmainfont{Times New Roman}
\newfontfamily\cyrillicfonttt{Courier New}
\usepackage[14pt]{extsizes}
\usepackage[left=20mm, top=20mm, right=20mm, bottom=20mm, footskip=10mm]{geometry}
\usepackage{ragged2e}
\usepackage{indentfirst}
\usepackage{graphicx}
\usepackage{caption}
\usepackage{xcolor}
\usepackage{listings}
\usepackage{booktabs}
\usepackage{hyperref}
\usepackage{amsmath}
\usepackage{amssymb}
\usepackage{array}
\usepackage{multirow}

\definecolor{link_color}{HTML}{0000EE}
\hypersetup{colorlinks=true,linkcolor=link_color,urlcolor=link_color,citecolor=link_color}

\lstset{
    language=Python,
    basicstyle=\cyrillicfonttt\small,
    numbers=left,
    numberstyle=\tiny\color{gray},
    keywordstyle=\color{blue}\bfseries,
    commentstyle=\color{gray}\itshape,
    stringstyle=\color{red},
    frame=single,
    captionpos=b,
    breaklines=true,
    breakatwhitespace=false,
    showstringspaces=false,
    keepspaces=true,
    columns=flexible,
    fontadjust=true,
    basewidth={0.5em,0.45em},
    tabsize=4,
}

\begin{document}

% ============================================================
% Титульный лист
% ============================================================
{
\thispagestyle{empty}
\begin{center}
    \hfill\break\hfill\break
    МИНИСТЕРСТВО НАУКИ И ВЫСШЕГО ОБРАЗОВАНИЯ РОССИЙСКОЙ ФЕДЕРАЦИИ\\[5pt]
    Федеральное государственное автономное образовательное учреждение высшего образования
    «Санкт-Петербургский политехнический университет Петра Великого»\\[10pt]
    Институт компьютерных наук и кибербезопасности\\[5pt]
    Направление: 02.03.01 Математика и компьютерные науки\\[30pt]
    {\large Отчет по дисциплине: «Вычислительная математика»}\\[30pt]
    \LARGE\bf «Лабораторная работа №4» (вариант 56).
\end{center}
\vspace{110pt}
\begin{tabular*}{460pt}{@{\extracolsep{\fill}} l r}
    Студент,\tabularnewline группы 5130201/40003 \hspace{50pt} & \line(1,0){100}\quad Четвергов И.\,С.\\
    & \\
    Руководитель,\tabularnewline Преподаватель & \line(1,0){121.5}\quad Пак В.\,Г.\\
    & \\ & \\ & \\
    & \guillemotleft\line(1,0){30}\guillemotright\line(1,0){100}\,2026\,г.
\end{tabular*}\\
\vfill
\begin{center}Санкт-Петербург, 2026\end{center}
\newpage
}

\tableofcontents
\newpage

% ============================================================
\section{Постановка задачи}
% ============================================================

Дан определённый интеграл (вариант 56):
\[
I = \int_{2}^{10} \frac{1}{x\ln(x)}\,dx.
\]

\textbf{Задание 1.} Вычислить приближённо интеграл по простейшим формулам
(левых прямоугольников, правых прямоугольников, центральных прямоугольников,
трапеций, парабол) при числе подынтервалов $n = 4$. Оценить погрешности.

\textbf{Задание 2.} Вычислить приближённо интеграл по формуле Ньютона-Котеса
с $n = 5$ узлами.

\textbf{Задание 3.} Вычислить приближённо интеграл по формуле Гаусса с $n = 4$
узлами. Оценить погрешность.

\subsection*{Точное значение интеграла}

Функция $f(x) = \dfrac{1}{x\ln(x)}$. Первообразная:
\[
F(x) = \ln|\ln(x)| + C.
\]
Точное значение интеграла:
\[
I = F(10) - F(2) = \ln(\ln 10) - \ln(\ln 2) = \ln\left(\frac{\ln 10}{\ln 2}\right)
  = \ln(3{,}321928) \approx 1{,}200973.
\]

\newpage

% ============================================================
\section{Задание 1. Простейшие квадратурные формулы}
% ============================================================

\subsection{Теоретические сведения}

Пусть отрезок $[a,b]$ разбит на $n$ равных частей с шагом $h = (b-a)/n$.
Узлы разбиения: $x_i = a + ih$, $i = 0, 1, \ldots, n$.

Формулы:
\[
I_L = h\sum_{i=0}^{n-1}f(x_i), \qquad
I_R = h\sum_{i=1}^{n}f(x_i), \qquad
I_C = h\sum_{i=0}^{n-1}f\!\left(x_i + \tfrac{h}{2}\right),
\]
\[
I_T = h\!\left(\frac{f(a)+f(b)}{2} + \sum_{i=1}^{n-1}f(x_i)\right)
\]
\[
I_S = \frac{h}{3}\!\left(f(x_0) + 4\!\sum_{\text{нечётн.}}\!f(x_i)
      + 2\!\sum_{\text{чётн.}}\!f(x_i) + f(x_n)\right).
\]

\subsection{Разбиение отрезка}

$n = 4$, $a = 2$, $b = 10$:
\[
h = \frac{b-a}{n} = \frac{10-2}{4} = 2.
\]

Узлы разбиения: $x_0 = 2{,}000$, $x_1 = 4{,}000$, $x_2 = 6{,}000$,
$x_3 = 8{,}000$, $x_4 = 10{,}000$.

\begin{table}[h!]
\centering
\caption{Значения функции в узлах разбиения}
\label{tab:s1_nodes}
\begin{tabular}{|c|c|l|}
\hline
$i$ & $x_i$ & $f(x_i) = \dfrac{1}{x_i\ln(x_i)}$ \\
\hline
0 & 2,0000 & $\dfrac{1}{2\ln 2} = \dfrac{1}{2 \cdot 0{,}693147} = 0{,}721348$ \\
1 & 4,0000 & $\dfrac{1}{4\ln 4} = \dfrac{1}{4 \cdot 1{,}386294} = 0{,}180337$ \\
2 & 6,0000 & $\dfrac{1}{6\ln 6} = \dfrac{1}{6 \cdot 1{,}791759} = 0{,}092733$ \\
3 & 8,0000 & $\dfrac{1}{8\ln 8} = \dfrac{1}{8 \cdot 2{,}079442} = 0{,}060166$ \\
4 & 10,0000 & $\dfrac{1}{10\ln 10} = \dfrac{1}{10 \cdot 2{,}302585} = 0{,}043429$ \\
\hline
\end{tabular}
\end{table}

\subsection{Формула левых прямоугольников}

\[
I_L = h\bigl(f(x_0)+f(x_1)+f(x_2)+f(x_3)\bigr)
\]
\[
= 2 \cdot\bigl(0{,}721348+0{,}180337+0{,}092733+0{,}060166\bigr)
\]
\[
= 2 \cdot 1{,}054584 = 2{,}109168.
\]
Фактическая погрешность:
\[
\Delta I_L = |I - I_L| = |1{,}200973 - 2{,}109168| = 0{,}908195.
\]

\subsection{Формула правых прямоугольников}

\[
I_R = h\bigl(f(x_1)+f(x_2)+f(x_3)+f(x_4)\bigr)
\]
\[
= 2 \cdot\bigl(0{,}180337+0{,}092733+0{,}060166+0{,}043429\bigr)
\]
\[
= 2 \cdot 0{,}376665 = 0{,}753330.
\]
Фактическая погрешность:
\[
\Delta I_R = |1{,}200973 - 0{,}753330| = 0{,}447643.
\]

\subsection{Формула центральных прямоугольников}

Середины подынтервалов: $\bar x_0 = 3{,}0$, $\bar x_1 = 5{,}0$,
$\bar x_2 = 7{,}0$, $\bar x_3 = 9{,}0$.

\begin{table}[h!]
\centering
\caption{Значения функции в серединах подынтервалов}
\label{tab:s1_mid}
\begin{tabular}{|c|c|c|}
\hline
$i$ & $\bar x_i$ & $f(\bar x_i)$ \\
\hline
0 & 3,0 & $\dfrac{1}{3\ln 3} = 0{,}303014$ \\
1 & 5,0 & $\dfrac{1}{5\ln 5} = 0{,}123552$ \\
2 & 7,0 & $\dfrac{1}{7\ln 7} = 0{,}076038$ \\
3 & 9,0 & $\dfrac{1}{9\ln 9} = 0{,}050313$ \\
\hline
\end{tabular}
\end{table}

\[
I_C = h\bigl(f(\bar x_0)+f(\bar x_1)+f(\bar x_2)+f(\bar x_3)\bigr)
\]
\[
= 2 \cdot\bigl(0{,}303014+0{,}123552+0{,}076038+0{,}050313\bigr)
\]
\[
= 2 \cdot 0{,}552917 = 1{,}105834.
\]
Фактическая погрешность:
\[
\Delta I_C = |1{,}200973 - 1{,}105834| = 0{,}095139.
\]

\subsection{Формула трапеций}

\[
I_T = h\!\left(\frac{f(x_0)+f(x_4)}{2}+f(x_1)+f(x_2)+f(x_3)\right)
\]
\[
= 2\!\left(\frac{0{,}721348+0{,}043429}{2}+0{,}180337+0{,}092733+0{,}060166\right)
\]
\[
= 2\cdot\bigl(0{,}382389 + 0{,}333236\bigr)
\]
\[
= 2 \cdot 0{,}715625 = 1{,}431250.
\]
Фактическая погрешность:
\[
\Delta I_T = |1{,}200973 - 1{,}431250| = 0{,}230277.
\]

\subsection{Формула парабол (Симпсона)}

При $n=4$: нечётные узлы --- $x_1$, $x_3$; чётный внутренний --- $x_2$.
\[
I_S = \frac{h}{3}\bigl(f(x_0)+4f(x_1)+2f(x_2)+4f(x_3)+f(x_4)\bigr)
\]
\[
= \frac{2}{3}\bigl(0{,}721348+4\cdot0{,}180337+2\cdot0{,}092733+4\cdot0{,}060166+0{,}043429\bigr)
\]
\[
= \frac{2}{3}\bigl(0{,}721348+0{,}721348+0{,}185466+0{,}240664+0{,}043429\bigr)
\]
\[
= \frac{2}{3}\cdot 1{,}912255 = 1{,}274837.
\]
Фактическая погрешность:
\[
\Delta I_S = |1{,}200973 - 1{,}274837| = 0{,}073864.
\]

\subsection{Оценка погрешностей}

Вычислим производные $f(x) = \dfrac{1}{x\ln(x)}$:
\[
f'(x) = -\frac{1}{x^2\ln(x)} - \frac{1}{x^2\ln^2(x)} = -\frac{\ln(x)+1}{x^2\ln^2(x)},
\]
\[
f''(x) = -\frac{d}{dx}\left(\frac{\ln(x)+1}{x^2\ln^2(x)}\right).
\]

После вычисления (через правило произведения и частного):
\[
f''(x) = -\frac{2\ln^3(x)+3\ln^2(x)+2\ln(x)+2}{x^3\ln^3(x)}.
\]

На $[2; 10]$ максимальные значения по модулю:
\[
M_1 = \max_{[2;10]}|f'(x)| = |f'(2)| = \frac{\ln(2)+1}{4\ln^2(2)} \approx 0{,}653855,
\]
\[
M_2 = \max_{[2;10]}|f''(x)| = |f''(2)| \approx 0{,}631478.
\]

Теоретические оценки погрешностей при $h = 2$, $b - a = 8$:

\begin{table}[h!]
\centering
\caption{Сравнение теоретических и фактических погрешностей}
\label{tab:s1_err}
\renewcommand{\arraystretch}{1.4}
\small
\begin{tabular}{|l|c|c|c|}
\hline
Формула & Оценка & Теор. оценка & Факт. погрешность \\
\hline
Лев./пр. прямоугольники
  & $\dfrac{M_1(b-a)h}{2}$
  & $\dfrac{0{,}653855\cdot 8\cdot 2}{2} = 5{,}230840$
  & $0{,}9082\;/\;0{,}4476$ \\[8pt]
\hline
Центр. прямоугольники
  & $\dfrac{M_2(b-a)h^2}{24}$
  & $\dfrac{0{,}631478\cdot 8\cdot 4}{24} = 0{,}841971$
  & $0{,}0951$ \\[8pt]
\hline
Трапеции
  & $\dfrac{M_2(b-a)h^2}{12}$
  & $\dfrac{0{,}631478\cdot 8\cdot 4}{12} = 1{,}683942$
  & $0{,}2303$ \\[8pt]
\hline
Парабол (Симпсон)
  & $\dfrac{M_4(b-a)h^4}{180}$
  & $\dfrac{1{,}265648\cdot 8\cdot 16}{180} = 1{,}789198$
  & $0{,}0739$ \\[8pt]
\hline
\end{tabular}
\end{table}

Здесь $M_4 = \max_{[2;10]}|f^{(4)}(x)| \approx 1{,}265648$. 
Фактические погрешности меньше или сравнимы с теоретическими оценками.
Наиболее точна формула центральных прямоугольников.

\newpage

% ============================================================
\section{Задание 2. Формула Ньютона-Котеса с $n = 5$ узлами}
% ============================================================

\subsection{Теоретические сведения}

Формула Ньютона-Котеса с $n$ узлами:
\[
\int_a^b f(x)\,dx \approx (b-a)\sum_{k=0}^{n-1} C_k^{(n)}\,f(x_k),
\]
где $x_k = a + kh$, $h = (b-a)/(n-1)$, коэффициенты $C_k^{(n)}$ берутся из таблицы.

При $n = 5$ узлах применяется \textbf{формула Буля}:
\[
\int_a^b f(x)\,dx \approx \frac{b-a}{90}\bigl[7f(x_0)+32f(x_1)+12f(x_2)+32f(x_3)+7f(x_4)\bigr].
\]
Остаточный член: $\displaystyle R^{(5)} = -\frac{8(b-a)}{945}\,h^6\,f^{(6)}(\xi)$,
$\xi\in[a,b]$.

\subsection{Вычисление}

Шаг:
\[
h = \frac{b-a}{n-1} = \frac{10-2}{4} = 2.
\]

Узлы интерполяции совпадают с табл.~\ref{tab:s1_nodes}:
\[
x_0 = 2{,}000,\quad x_1 = 4{,}000,\quad x_2 = 6{,}000,
\quad x_3 = 8{,}000,\quad x_4 = 10{,}000.
\]

Значения функции из табл.~\ref{tab:s1_nodes}. Подставляем в формулу Буля:
\begin{align*}
&7f(x_0)+32f(x_1)+12f(x_2)+32f(x_3)+7f(x_4)\\
&= 7\cdot0{,}721348 + 32\cdot0{,}180337 + 12\cdot0{,}092733
   + 32\cdot0{,}060166 + 7\cdot0{,}043429\\
&= 5{,}049436 + 5{,}770784 + 1{,}112796 + 1{,}925312 + 0{,}304003\\
 = 14{,}162331.
\end{align*}

\[
\boxed{I^*_\text{НК} = \frac{8}{90}\cdot 14{,}162331
= 0{,}0\overline{8}\cdot 14{,}162331 \approx 1{,}258036.}
\]

\subsection{Погрешность}

Фактическая погрешность:
\[
\Delta I^*_\text{НК} = |I - I^*_\text{НК}| = |1{,}200973 - 1{,}258036| = 0{,}057063.
\]

Теоретическая оценка остаточного члена. Шестая производная $f^{(6)}(x)$
вычислена численно; на $[2; 10]$:
\[
M_6 = \max_{[2;10]}\bigl|f^{(6)}(x)\bigr| \approx 2{,}458637 \quad (\text{при }x=2).
\]

\[
|R^{(5)}| \leq \frac{8(b-a)}{945}\,h^6\,M_6
= \frac{8\cdot 8}{945}\cdot(2)^6\cdot 2{,}458637\]
\[= 0{,}067749\cdot 64\cdot 2{,}458637 = 10{,}661189.
\]

Фактическая погрешность $0{,}057063 < 10{,}661189$ --- оценка выполнена (консервативна).

\subsection{Сравнение формул Ньютона-Котеса}

\begin{table}[h!]
\centering
\caption{Результаты формул Ньютона-Котеса}
\label{tab:nc_comp}
\begin{tabular}{|c|l|c|c|}
\hline
$n$ & Название & $I^*$ & $\Delta I$ \\
\hline
2 & Трапеция (весь отрезок) & 0,382389 & 0,818584 \\
3 & Симпсон 1/3              & 1,157292 & 0,043681 \\
4 & Три восьмых              & 1,218715 & 0,017742 \\
5 & Буля                     & 1,258036 & 0,057063 \\
\hline
\end{tabular}
\end{table}

С ростом числа узлов точность формул Ньютона-Котеса возрастает (за исключением формулы Буля с $n=5$, где функция имеет особенность).

\newpage

% ============================================================
\section{Задание 3. Формула Гаусса с $n = 4$ узлами}
% ============================================================

\subsection{Теоретические сведения}

Для вычисления интеграла $\displaystyle\int_a^b f(x)\,dx$ по формуле Гаусса с $n$
узлами на произвольном отрезке $[a,b]$ узлы пересчитываются из стандартных
узлов $t_i \in [-1,1]$ по формуле:
\[
x_i = \frac{a+b}{2} + \frac{b-a}{2}\,t_i,
\]
а приближённое значение интеграла:
\[
\int_a^b f(x)\,dx \approx \frac{b-a}{2}\sum_{i=1}^{n} A_i\,f(x_i),
\]
где $A_i$ --- веса формулы (не меняются при переходе к $[a,b]$).

\subsection{Узлы и веса для $n = 4$}

\begin{table}[h!]
\centering
\caption{Узлы $t_i$ и веса $A_i$ формулы Гаусса при $n=4$ на $[-1,1]$}
\label{tab:gauss_tab}
\begin{tabular}{|c|c|c|}
\hline
$i$ & $t_i$ & $A_i$ \\
\hline
1 & $-0{,}861136312$ & $0{,}347854845$ \\
2 & $-0{,}339981044$ & $0{,}652145155$ \\
3 & $+0{,}339981044$ & $0{,}652145155$ \\
4 & $+0{,}861136312$ & $0{,}347854845$ \\
\hline
\end{tabular}
\end{table}

\subsection{Пересчёт узлов на отрезок $[2; 10]$}

\[
\frac{a+b}{2} = \frac{2+10}{2} = 6, \qquad
\frac{b-a}{2} = \frac{10-2}{2} = 4.
\]
\[
x_i = 6 + 4\,t_i.
\]
\begin{align*}
x_1 &= 6 + 4\cdot(-0{,}861136) = 6 - 3{,}444544 = 2{,}555456,\\
x_2 &= 6 + 4\cdot(-0{,}339981) = 6 - 1{,}359924 = 4{,}640076,\\
x_3 &= 6 + 4\cdot(+0{,}339981) = 6 + 1{,}359924 = 7{,}359924,\\
x_4 &= 6 + 4\cdot(+0{,}861136) = 6 + 3{,}444544 = 9{,}444544.
\end{align*}

\subsection{Вычисление значений функции}
Показаны в таблице~\ref{tab:gauss_calc}.

\begin{table}[h!]
\centering
\caption{Узлы, веса и значения функции для формулы Гаусса ($n=4$) на $[2; 10]$}
\label{tab:gauss_calc}
\renewcommand{\arraystretch}{1.3}
\begin{tabular}{|c|c|c|c|c|c|}
\hline
$i$ & $t_i$ & $x_i$ & $A_i$ & $f(x_i)$ & $A_i\,f(x_i)$ \\
\hline
1 & $-0{,}861136312$ & $2{,}555456$ & $0{,}347854845$ & $0{,}569276$ & $0{,}197940$ \\
2 & $-0{,}339981044$ & $4{,}640076$ & $0{,}652145155$ & $0{,}158852$ & $0{,}103537$ \\
3 & $+0{,}339981044$ & $7{,}359924$ & $0{,}652145155$ & $0{,}077752$ & $0{,}050699$ \\
4 & $+0{,}861136312$ & $9{,}444544$ & $0{,}347854845$ & $0{,}046386$ & $0{,}016131$ \\
\hline
\multicolumn{5}{|r|}{$\displaystyle\sum_{i=1}^{4} A_i\,f(x_i)$} & $0{,}368307$ \\
\hline
\end{tabular}
\end{table}

\subsection{Вычисление интеграла}

\[
I^*_\text{Г} = \frac{b-a}{2}\sum_{i=1}^{4}A_i\,f(x_i)
= 4\cdot\bigl(0{,}197940+0{,}103537+0{,}050699+0{,}016131\bigr)
\]
\[
= 4\cdot 0{,}368307 = 1{,}473228.
\]

\subsection{Оценка погрешности}

Точное значение: $I \approx 1{,}200973$. Фактическая погрешность:
\[
\Delta I^*_\text{Г} = |I - I^*_\text{Г}| = |1{,}200973 - 1{,}473228| = 0{,}272255.
\]

Формула Гаусса при $n=4$ имеет алгебраическую степень точности $2n-1=7$,
т.\,е. точна для всех многочленов степени не выше~7. Однако функция
$f(x) = \dfrac{1}{x\ln x}$ имеет особенность при $x \to 1^+$ и быстро
убывает, что требует более высокого порядка квадратуры для хорошей точности.
Погрешность сравнительно велика из-за характера функции.

\newpage

% ============================================================
\section{Выводы}
% ============================================================

В лабораторной работе вычислен интеграл
$\displaystyle\int_2^{10}\frac{1}{x\ln(x)}\,dx \approx 1{,}200973$
тремя группами методов.

\begin{enumerate}
\item \textbf{Простейшие формулы} ($n=4$ подынтервала): формулы левых и правых
прямоугольников дают грубые приближения (погрешность $\sim 0{,}45$--$0{,}91$),
так как функция быстро убывает на отрезке. Центральные прямоугольники дают
приемлемый результат (погрешность $0{,}0951$). Формула трапеций даёт
погрешность $0{,}2303$, а Симпсона --- $0{,}0739$.

\item \textbf{Формула Ньютона-Котеса с $n=5$ узлами} (Буля) даёт погрешность
$0{,}057063$, что лучше формулы Симпсона. Формулы Ньютона-Котеса
хорошо работают для гладких функций.

\item \textbf{Формула Гаусса с $n=4$ узлами} показывает худший результат
($\Delta I^*_\text{Г} \approx 0{,}272255$) по сравнению с другими методами.
Это объясняется характером функции (логарифмическая особенность)
и тем, что Гаусс оптимален для многочленов, а не для таких функций.
\end{enumerate}

\newpage

% ============================================================
\section*{Приложение}
\addcontentsline{toc}{section}{Приложение}

\begin{lstlisting}[caption={Численное интегрирование (Python)}, label=lst:code]
import numpy as np
from scipy import integrate

f = lambda x: 1.0 / (x * np.log(x))
a, b = 2.0, 10.0

# Tochnoe znachenie
I_exact, _ = integrate.quad(f, a, b)
print(f"Tochnoe: {I_exact:.6f}")

# ===== Zadanie 1: prosteyshie formuly, n=4 =====
n = 4
h = (b - a) / n         # 2.0
xi = [a + i*h for i in range(n+1)]

I_L = h * sum(f(xi[i]) for i in range(n))
I_R = h * sum(f(xi[i]) for i in range(1, n+1))
I_C = h * sum(f(a + (i+0.5)*h) for i in range(n))
I_T = h * (f(a)/2 + sum(f(xi[i]) for i in range(1,n)) + f(b)/2)
I_S = h/3 * (f(a) + 4*sum(f(xi[i]) for i in range(1,n,2))
           + 2*sum(f(xi[i]) for i in range(2,n,2)) + f(b))

print(f"Lev. pr.: {I_L:.6f}  err={abs(I_exact-I_L):.6f}")
print(f"Pr.  pr.: {I_R:.6f}  err={abs(I_exact-I_R):.6f}")
print(f"Centr.  : {I_C:.6f}  err={abs(I_exact-I_C):.6f}")
print(f"Trapec. : {I_T:.6f}  err={abs(I_exact-I_T):.6f}")
print(f"Simpson : {I_S:.6f}  err={abs(I_exact-I_S):.6f}")

# ===== Zadanie 2: Newton-Kotes n=5 (formula Bulya) =====
h_nc = (b - a) / 4      # 2.0
xnc  = [a + i*h_nc for i in range(5)]
fv   = [f(x) for x in xnc]
I_NC = (b-a)/90 * (7*fv[0]+32*fv[1]+12*fv[2]+32*fv[3]+7*fv[4])
print(f"NK (n=5): {I_NC:.6f}  err={abs(I_exact-I_NC):.6f}")

# ===== Zadanie 3: Gauss n=4 =====
t = [-0.861136312, -0.339981044, 0.339981044, 0.861136312]
A = [ 0.347854845,  0.652145155, 0.652145155, 0.347854845]
mid  = (a + b) / 2      # 6.0
half = (b - a) / 2      # 4.0
x_g  = [mid + half*ti for ti in t]
I_G  = half * sum(A[i]*f(x_g[i]) for i in range(4))
print(f"Gauss(n=4): {I_G:.6f}  err={abs(I_exact-I_G):.6f}")
\end{lstlisting}

\end{document}